\documentclass[a4paper,11pt]{article}
\usepackage{amsmath,amsthm,amsfonts,amssymb,amscd,amstext,vmargin,graphics,graphicx,tabularx,multicol} 
\usepackage[francais]{babel}
\usepackage[utf8]{inputenc}  
\usepackage[T1]{fontenc} 
\usepackage{pstricks-add,tikz,tkz-tab,variations}
\usepackage[autolanguage,np]{numprint} 
\usepackage{calc}

\setmarginsrb{1.5cm}{0.5cm}{1cm}{0.5cm}{0cm}{0cm}{0cm}{0cm} %Gauche, haut, droite, haut
\newcounter{numexo}
\newcommand{\exo}[1]{\stepcounter{numexo}\noindent{\bf Exercice~\thenumexo} : }
\reversemarginpar

\newcommand{\bmul}[1]{\begin{multicols}{#1}}
\newcommand{\emul}{\end{multicols}}

\newcounter{enumtabi}
\newcounter{enumtaba}
\newcommand{\q}{\stepcounter{enumtabi} \theenumtabi.  }
\newcommand{\qa}{\stepcounter{enumtaba} (\alph{enumtaba}) }
\newcommand{\initq}{\setcounter{enumtabi}{0}}
\newcommand{\initqa}{\setcounter{enumtaba}{0}}

\newcommand{\be}{\begin{enumerate}}
\newcommand{\ee}{\end{enumerate}}
\newcommand{\bi}{\begin{itemize}}
\newcommand{\ei}{\end{itemize}}
\newcommand{\bp}{\begin{pspicture*}}
\newcommand{\ep}{\end{pspicture*}}
\newcommand{\bt}{\begin{tabular}}
\newcommand{\et}{\end{tabular}}
\renewcommand{\tabularxcolumn}[1]{>{\centering}m{#1}} %(colonne m{} centrée, au lieu de p par défault) 
\newcommand{\tnl}{\tabularnewline}

\newcommand{\trait}{\noindent \rule{\linewidth}{0.2mm}}
\newcommand{\hs}[1]{\hspace{#1}}
\newcommand{\vs}[1]{\vspace{#1}}

\newcommand{\N}{\mathbb{N}}
\newcommand{\Z}{\mathbb{Z}}
\newcommand{\R}{\mathbb{R}}
\newcommand{\C}{\mathbb{C}}
\newcommand{\Dcal}{\mathcal{D}}
\newcommand{\Ccal}{\mathcal{C}}
\newcommand{\mc}{\mathcal}

\newcommand{\vect}[1]{\overrightarrow{#1}}
\newcommand{\ds}{\displaystyle}
\newcommand{\eq}{\quad \Leftrightarrow \quad}
\newcommand{\vecti}{\vec{\imath}}
\newcommand{\vectj}{\vec{\jmath}}
\newcommand{\Oij}{(O;\vec{\imath}, \vec{\jmath})}
\newcommand{\OIJ}{(O;I,J)}


\newcommand{\reponse}[1][1]{%
\multido{}{#1}{\makebox[\linewidth]{\rule[0pt]{0pt}{20pt}\dotfill}
}}

\newcommand{\titre}[5] 
% #1: titre #2: haut gauche #3: bas gauche #4: haut droite #5: bas droite
{
\noindent #2 \hfill #4 \\
#3 \hfill #5

\vspace{-1.6cm}

\begin{center}\rule{6cm}{0.5mm}\end{center}
\vspace{0.2cm}
\begin{center}{\large{\textbf{#1}}}\end{center}
\begin{center}\rule{6cm}{0.5mm}\end{center}
}



\begin{document}
\pagestyle{empty}
\titre{Séance d'AP  : les puissances positives de 10}{}{}{4ème}{}

\vspace*{0.2cm}

\setlength{\fboxrule}{2pt}
\begin{flushleft}
\framebox{\begin{minipage}{\linewidth}

\vspace*{0.5cm}

\underline{\textbf{{\large Cours}}}\\

\textbf{Les puissances de 10} servent souvent à écrire les nombres très grands ou très petits.\\




\textbf{Définition 1:} Un multiple de 10 peut s'écrire sous la forme d'un produit $10^{n}$, où est un entier positif non nul.
 \begin{center}
  $10^{n}$ est le produit de $n$ facteurs égaux à 10 : \hspace*{1.5cm}  $10^{n} = \underbrace{10 \times ... \times 10}_{n fois}= \underbrace{10...0}_{n zero}$ 
   \end{center}  

\textbf{Remarques :} \hspace*{1cm}  Si n=1 alors $10^{1}=10$\\
\hspace*{3.5cm}  Si n=2 alors $10^{2}=10 \times 10 = 100$\\
\hspace*{3.5cm}  Si n=3 alors $10^{3}=10 \times 10 \times 10 = 1000$\\
\hspace*{3.5cm}  Si n=4 alors $10^{4}=10 \times 10 \times 10 \times 10= 10000$\\

\textbf{Applications : } 1 Méga octets = 1 millions d'octets = $10^{6}  $ o = 1 MO
\end{minipage}}
\end{flushleft}


\renewcommand{\arraystretch}{1.5}


\exo \\
Un nouveau virus informatique est détecté ; il se nomme « Puis Sens Disse », nom de code : PS10. \\
Ce virus envahit un secteur du disque dur et le rend inutilisable. Au bout d’une heure, PS10 a infecté dix nouveaux
secteurs, et ainsi de suite : durant chaque heure suivante, chaque secteur contamine dix nouveaux secteurs sains.\\

\q Combien de secteur sont contaminés :\\

\includegraphics[scale=1]{virus.eps} \\
\q L’ordinateur est contaminé depuis 8 heures. Sur l’écran, on peut lire : "$10^{8}$ secteurs détruits durant la 8ème heure". Pourquoi a-t-on cette notation $10^{8}$? Donner l’écriture de ce nombre.\\
\reponse[2]\\
\q Cent mille millions de secteurs viennent d’être contaminés. Quel message peut-on lire à l’écran ? \\
\reponse[1]\\

\vspace*{1cm}


\bmul{2}
\exo \\
Écrire les nombres suivants sous forme de puissance de 10. \\



\noindent 1 000 = . . . . . . . . .\\
100 000 000 = . . . . . . . \\
100 = . . . . . . . . \\
10 000 = . . . . . . . . . \\
10 000 000 000 = . . . . . . . . . . .\\
 

\columnbreak

\exo \\
Calculer les expressions suivantes.

\bmul{2}

$2,345 \times 10^{2} = $ . . . . .\\


$0,03 \times 10^{2} = $ . . . . .\\

$752 \times 10^{1} = $ . . . . .\\

\columnbreak


$35,98  \times 10^{5} = $ . . . . .\\

$0,89 \times 10^{7} = $ . . . . .\\

$3,592  \times 10^{4} = $ . . . . .\\

\emul




\emul

\newpage

\vspace*{0.2cm}

\setlength{\fboxrule}{2pt}
\begin{flushleft}
\framebox{\begin{minipage}{\linewidth}


\underline{\textbf{Cours}}\\


\textbf{Définition 2 :}\textbf{ L'écriture scientifique }d'un nombre décimal est l'unique écriture de ce nombre sous la forme $a \times 10^{n}$, où:
\bi
\item $a$ est un nombre décimal compris entre 1 et 10, 10 étant exclu ;
\item $n$ est un nombre entier relatif (positif ou négatif).\\
\ei
\vspace*{0.2cm}

\textbf{Exemples :}  $2,145 \times 10^{5}$ est l'écriture scientifique du nombre 214 500\\
 \hspace*{ 1.8cm}               $1,25 \times 10^{-2}$  est l'écriture scientifique du nombre 0, 0125


\vspace*{0.5cm}
\end{minipage}}
\end{flushleft}

\vspace*{0.25cm}
\bmul{2}
\exo \\
Associer les écritures scientifiques avec les nombres auxquels elles correspondent. \\



\noindent $2,718 \times 10^{3}$ \hspace*{2cm} 2,718\\
$2,718 \times  10^{0}$    \hspace*{2cm} 2 718\\
$2,718 \times  10^{2}$    \hspace*{2cm} 27,18\\
$2,718 \times  10^{1}$    \hspace*{2cm} 27 180\\
$2,718 \times  10^{5}$   \hspace*{2cm} 271,8\\
$2,718 \times  10^{-2}$   \hspace*{1.8cm} 0,02718\\


\columnbreak

\exo \\
Comparer ces nombres en écriture scientifique, avec les signes $<$ ou $>$:\\

\initqa 
\qa $9,45 \times 10^{8} $ . . . . $8,31 \times 10^{8} $\\

\qa $9 \times 10^{3} $ . . . . $9 \times 10^{2} $\\

\qa $3,5 \times 10^{13} $ . . . . $2,65 \times 10^{12} $\\

\qa $7,2 \times 10^{15} $ . . . . $7,2 \times 10^{13} $\\

\qa $1,6 \times 10^{9} $ . . . . $1,5 \times 10^{10} $\\

\emul

\vspace*{0.25cm}

\exo \\ Le tableau suivant donne les distances, en kilomètres, entre les planètes du système solaire et le Soleil :\\


\begin{tabular}{|c|c|c|c|}
\hline 
\textbf{Planètes} & \textbf{Distances} & \textbf{Écritures scientifiques }& \textbf{Ordre des planètes} \\ 
\hline 
\textbf{Saturne}  & $1,427 \times 10^{9}$ &  &  \\ 
\hline 
\textbf{Terre} &  $149,6 \times 10^{6}$  &  &  \\ 
\hline 
\textbf{Jupiter} &  $7,78 \times 10^{8}$  &  &  \\ 
\hline 
\textbf{Neptune} &  $44,97 \times 10^{8}$  &  &  \\ 
\hline 
\textbf{Mercure} &  $57,9 \times 10^{6}$  &  & 1 \\ 
\hline 
\textbf{Vénus} &  $108 \times 10^{6}$  &  &  \\ 
\hline 
\textbf{Mars} &  $228 \times 10^{6}$  &  &  \\ 
\hline 
\textbf{Uranus} &  $2,871 \times 10^{9}$  &  &  \\ 
\hline 
\end{tabular} 

\vspace*{0.25cm}

$\rightarrow $ Après avoir écrit les distances des planètes au Soleil en écriture scientifique,indique leur ordre \textbf{de la plus proche à la plus éloignée du soleil}.





\setlength{\fboxrule}{2pt}
\begin{flushleft}
\framebox{\begin{minipage}{\linewidth}

\vspace*{0.5cm}

\begin{center}
\underline{\textbf{Pour information}}
\end{center}


\begin{tabular}{|c|c|c|c|c|c|}
\hline 
Préfixe & Symbole & Puissances de 10 & Préfixe & Symbole & Puissance de 10 \\ 
\hline 
Téra & T & $10^{12}$ & Kilo & k & $10^{3}$\\ 
\hline 
Giga & G & $10^{9}$ & Hecto & h & $10^{2}$ \\ 
\hline 
Méga & M & $10^{6}$ & Déca & da & $10^{1}$ \\ 
\hline 
\end{tabular} 


\vspace*{0.5cm}
\end{minipage}}
\end{flushleft}

\newpage


\end{document}
